\documentclass[runningheads]{llncs}
%
\usepackage[T1]{fontenc}
% T1 fonts will be used to generate the final print and online PDFs,
% so please use T1 fonts in your manuscript whenever possible.
% Other font encondings may result in incorrect characters.
%
\usepackage{graphicx}
% Used for displaying a sample figure. If possible, figure files should
% be included in EPS format.
%
% If you use the hyperref package, please uncomment the following two lines
% to display URLs in blue roman font according to Springer's eBook style:
%\usepackage{color}
%\renewcommand\UrlFont{\color{blue}\rmfamily}
%\urlstyle{rm}

\usepackage{amsmath}
\usepackage{tikz}
\usetikzlibrary{arrows.meta,positioning,fit,calc}

\usepackage{caption}
\usepackage{subcaption}

\begin{document}
%
\title{Pattern Layout Optimization for \\Seam-Aware Texture Alignment}

\section{Formal Problem Specification: Simultaneous Topology and Layout Optimization}
\label{chap:formal_spec}

\subsection{Problem Domain and Inputs}

We define the problem as a tuple $\Psi = \langle M, \mathcal{T}, \mathcal{L}, \Omega \rangle$, where:

\begin{itemize}
    \item $M = (V, E, F)$ is the input discrete manifold (triangle mesh) representing the 3D garment surface.
    \item $\mathcal{T}: \mathbf{R}^2 \to \mathbf{R}^k$ is the periodic texture field defined on the fabric domain. We assume $\mathcal{T}$ is characterized by a lattice basis $B = \{\mathbf{u}, \mathbf{v}\}$ and a signal function $\tau(\mathbf{x})$.
    \item $\mathcal{L} = \{L_1, \dots, L_m\}$ is a set of anatomical landmarks on $M$, where each $L_j \subset F$ is a constrained region (subset of faces) defining the valid domain for a specific seam keypoint.
    \item $\Omega = [0, W] \times [0, \infty)$ is the fabric domain, a semi-infinite strip of fixed width $W \in \mathbf{R}^+$.
\end{itemize}

\subsection{Decision Variables (Genotype)}

A candidate solution is a vector $\mathbf{x} = \langle \boldsymbol{\delta}, \boldsymbol{\rho}, \boldsymbol{\alpha}, \pi, h \rangle$ defined as follows:

\subsubsection{Geometric Topology Parameters ($\boldsymbol{\delta}$)}
Let $K = \{k_1, \dots, k_m\}$ be the set of mobile seam endpoints.
$\boldsymbol{\delta} \in [0, 1]^m$ is a vector of normalized coordinates.
For each $j \in \{1, \dots, m\}$, the position of keypoint $k_j$ is given by a mapping function $\Gamma_j: [0, 1] \to M$, where $\Gamma_j(\delta_j)$ returns a point on the surface within the region $L_j$.
The set of seams $\mathcal{S}$ is the fixed topological graph connecting vertices in $K$. A specific realization of seam geometry is the set of geodesic paths $G(\boldsymbol{\delta}) = \{g_s \mid s \in \mathcal{S}\}$, where $g_s$ is the shortest path on $M$ between the computed locations of its endpoints.

\subsubsection{Fabrication Orientation ($\boldsymbol{\rho}$)}
Let $P(\boldsymbol{\delta}) = \{p_1, \dots, p_n\}$ be the set of surface charts resulting from cutting $M$ along $G(\boldsymbol{\delta})$.
$\boldsymbol{\rho} \in \{0, 1, 2, 3\}^n$ is the discrete orientation vector.
For each chart $p_i$, the fabrication grain rotation is $\theta_i = \rho_i \cdot \frac{\pi}{2}$.

\subsubsection{Alignment Logic ($\boldsymbol{\alpha}$)}
Let $\mathcal{S}_{adj} \subset \mathcal{S}$ be the subset of internal seams connecting two distinct charts.
$\boldsymbol{\alpha} \in \{0, 1\}^{|\mathcal{S}_{adj}|}$ is the binary constraint vector.
For a seam $s \in \mathcal{S}_{adj}$ connecting charts $p_i$ and $p_j$:
\begin{itemize}
    \item If $\alpha_s = 0$, the seam is \textit{Relaxed}. Texture continuity is optimized but not enforced.
    \item If $\alpha_s = 1$, the seam is \textit{Locked}. The charts $p_i$ and $p_j$ are subject to a congruency constraint during flattening and a rigid coupling constraint during nesting.
\end{itemize}

\subsubsection{Nesting Logic ($\pi, h$)}
$\pi \in S_n$ is a permutation of $\{1, \dots, n\}$ defining the placement priority.
$h \in \mathcal{H}$ is an index selecting a heuristic placement function from the set $\mathcal{H} = \{ \text{LeftBottom}, \text{Skyline}, \text{LargestFirst} \}$.

\subsection{Phenotype Construction Pipeline}
The mapping $\Phi: \mathbf{x} \to L$ (from genotype to layout) proceeds in three deterministic stages.

\subsubsection{Stage 1: Constrained Flattening}
Given the cuts $G(\boldsymbol{\delta})$, we compute the 2D flattening $\phi_i: p_i \to \mathbf{R}^2$ for each chart. The flattening minimizes the isometric distortion energy $E_{iso}(\phi_i)$ subject to boundary conditions defined by $\boldsymbol{\alpha}$.

Let $\partial p_{i,s}$ denote the boundary curve of chart $p_i$ corresponding to seam $s$.
For every active constraint $\alpha_s = 1$ connecting $p_i, p_j$:
\begin{equation}
    \text{minimize } \sum E_{iso} \quad \text{s.t. } \quad \phi_i(\partial p_{i,s}) \cong \phi_j(\partial p_{j,s})
\end{equation}
where $\cong$ denotes geometric congruence up to a rigid transformation consistent with the texture symmetries of $\mathcal{T}$.
If $\alpha_s = 0$, the boundary is free.

\subsubsection{Stage 2: Cluster Assembly}
We form a set of rigid nesting items $\mathcal{I}$. Initially, $\mathcal{I} = \{ \phi_1(p_1), \dots, \phi_n(p_n) \}$.
We process constraints $\boldsymbol{\alpha}$ iteratively. For each $s$ with $\alpha_s = 1$:
\begin{enumerate}
    \item Identify items $I_u, I_v \in \mathcal{I}$ containing the adjacent boundaries of $s$.
    \item Compute the unique rigid transformation $T_{v \to u}$ that aligns the boundary $\partial p_{j,s}$ with $\partial p_{i,s}$ such that the texture lattice $B$ is continuous.
    \item Merge $I_u$ and $I_v$ into a new super-item $I_{new} = I_u \cup T_{v \to u}(I_v)$.
    \item Update $\mathcal{I} \leftarrow (\mathcal{I} \setminus \{I_u, I_v\}) \cup \{I_{new}\}$.
\end{enumerate}
The process terminates when all active constraints are satisfied.

\subsubsection{Stage 3: Heuristic Nesting}
Let $\mathcal{I}_{final}$ be the set of assembled clusters. The nesting algorithm places each item $I \in \mathcal{I}_{final}$ into $\Omega$.
\begin{enumerate}
    \item \textbf{Rotation:} Apply rotation $R(\theta_{base})$ to $I$. Note that for clusters, relative rotations of sub-parts are fixed in Stage 2; only the global cluster orientation varies.
    \item \textbf{Ordering:} Sort $\mathcal{I}_{final}$ according to the permutation $\pi$ (using the index of the primary root component of each cluster).
    \item \textbf{Placement:} For each $I_k$, find position $\mathbf{t}^* \in \Omega$ such that:
    \begin{equation}
        \mathbf{t}^* = \underset{\mathbf{t}}{\arg\min} \,\, \text{Cost}_h(\mathbf{t}, I_k, \mathcal{L}_{prev})
    \end{equation}
    subject to:
    \begin{align}
        (I_k + \mathbf{t}) \cap (I_{prev} + \mathbf{t}_{prev}) = \emptyset \quad \forall I_{prev} \in \mathcal{L}_{prev} \\
        (I_k + \mathbf{t}) \subset \Omega
    \end{align}
    where $\text{Cost}_h$ is the cost function of heuristic $h$ (e.g., max y-coordinate).
    \item \textbf{Lattice Snapping:} If $I_k$ contains no locked seams ($\alpha=0$ for all internal edges), update $\mathbf{t}^* \leftarrow \text{Snap}(\mathbf{t}^*, B)$ to align with the texture phase.
\end{enumerate}

\subsection{Objective Functions}

The fitness vector $\mathbf{F}(\mathbf{x}) = [f_1, f_2, f_3]^T$ is evaluated on the final layout.

\paragraph{1. Fabric Efficiency ($f_1$)}
\begin{equation}
    f_1(\mathbf{x}) = \max_{p \in \text{Layout}} (y_{max}(p))
\end{equation}
The total length of the fabric marker used.

\paragraph{2. Texture Discontinuity ($f_2$)}
Let $\mathcal{S}_{relaxed} = \{s \in \mathcal{S}_{adj} \mid \alpha_s = 0\}$.
\begin{equation}
    f_2(\mathbf{x}) = \sum_{s \in \mathcal{S}_{relaxed}} \int_{0}^{len(s)} \| \nabla \mathcal{T}(\mathbf{u}_i(l)) - \nabla \mathcal{T}(\mathbf{u}_j(l)) \|^2 \, dl
\end{equation}
The integral of texture signal deviation along relaxed seams. Note that for locked seams ($\alpha_s=1$), the mismatch is structurally zero by construction in Stage 2.

\paragraph{3. Geometric Distortion ($f_3$)}
\begin{equation}
    f_3(\mathbf{x}) = \sum_{i=1}^n E_{iso}(\phi_i)
\end{equation}
The sum of conformal energy from Stage 1. This captures the "stress" induced by enforcing straight/congruent seams on curved surfaces.

%\end{chapter}

\end{document}